\uuid{cP7w}
\exo7id{4871}
\titre{exo7 4871}
\auteur{quercia}
\organisation{exo7}
\datecreate{2010-03-17}
\isIndication{false}
\isCorrection{false}
\chapitre{Géométrie affine dans le plan et dans l'espace}
\sousChapitre{Applications affines}
\cpge{mpsi-mp2i.sous-espaces-affines, ptsi.espaces-vectoriels.b, mpsi-mp2i.applications-lineaires.c, mp-mpi.topologie-espaces-normes.j, pc.espaces-vectoriels-normes.g, psi.espaces-vectoriels-normes.g, bcpst2.valeurs-propres-vecteurs-propres.b}
\module{Géométrie}
\niveau{L2}
\difficulte{}

\contenu{
\question{
Soit ${\cal E}$ un espace affine de dimension finie et
$f : {\cal E} \to {\cal E}$ affine.
Montrer que $f$ admet un unique point fixe si et seulement si $1$ n'est pas
valeur propre de $\vec f$.
}
}
